% ============================================================================
% Section 3: Mathematical Formalism
% ============================================================================

\section{Mathematical formalism}
\label{sec:formalism}

This section develops the mathematical machinery of UHM from the axioms
stated in Section~\ref{sec:axioms}. Every definition and formula is derived
from the primitive $\mathfrak{T} = (\ShInf(\mathcal{C}),\; J_{\mathrm{Bures}},\; \omega_0)$;
nothing is postulated ad hoc. The complete proofs are available at
\url{https://holon.sh}; the key results are presented here with sufficient
detail for verification.

\subsection{The coherence matrix $\Gam$}
\label{subsec:gamma}

\begin{definition}[Coherence matrix]
The coherence matrix is an object of $\mathcal{C} = \mathbf{DensityMat}(\mathbb{C}^7)$:
\begin{equation}
  \Gam = \sum_{i,j \in \{A,S,D,L,E,O,U\}} \gamma_{ij}\, |i\rangle\langle j|
  \label{eq:gamma-def}
\end{equation}
satisfying three properties:
\begin{enumerate}[leftmargin=2em]
  \item \textbf{Hermiticity:} $\Gam^\dagger = \Gam$, i.e., $\gamma_{ij} = \gamma_{ji}^*$
  \item \textbf{Positive semi-definiteness:} $\langle\psi|\Gam|\psi\rangle \geq 0$
        for all $|\psi\rangle \in \Hil$, equivalently all eigenvalues $\lambda_k \geq 0$
  \item \textbf{Normalization:} $\Tr(\Gam) = 1$
\end{enumerate}
\end{definition}

\paragraph{Degrees of freedom.}
A $7 \times 7$ Hermitian matrix has $7^2 = 49$ real parameters. The trace constraint
removes one, leaving 48 independent parameters. The automorphism group $G_2 = \mathrm{Aut}(\Oct)$
acts as a gauge group with $\dim(G_2) = 14$, yielding $48 - 14 = 34$ physically
distinguishable (gauge-invariant) parameters.

\paragraph{Diagonal and off-diagonal.}
The diagonal elements $\gamma_{kk}$ represent the weight of each dimension (how much
of the system's ``resources'' are allocated to perception, structure, dynamics, etc.).
The off-diagonal elements $\gamma_{ij}$ ($i \neq j$) represent \emph{coherences} ---
quantum correlations between dimensions. The necessity of complex coherences
($\gamma_{ij} \in \mathbb{C}$) for non-trivial Gap structure is proved as T-132~\status{T}.

\subsection{Purity and the viability threshold}
\label{subsec:purity}

\begin{definition}[Purity]
\begin{equation}
  P(\Gam) := \Tr(\Gam^2) = \sum_i \gamma_{ii}^2 + \sum_{i \neq j} |\gamma_{ij}|^2
  \label{eq:purity}
\end{equation}
\end{definition}

Purity ranges from $P = 1/N = 1/7 \approx 0.143$ (maximally mixed state $\Gam = I/7$,
complete thermal equilibrium) to $P = 1$ (pure state, maximum organization). It measures
how structured the system is --- equivalently, how distinguishable from noise.

\begin{theorembox}[Critical purity (T-160) \status{T}]
For a system with $\Gam \in \Dens(\mathbb{C}^7)$, the viability threshold is:
\begin{equation}
  \Pcrit = \frac{2}{N} = \frac{2}{7} \approx 0.286
  \label{eq:pcrit}
\end{equation}
Below this threshold, the system is informationally indistinguishable from thermal noise
in the Bures metric and cannot sustain regeneration. Five independent derivations converge
on this value (Appendix~\ref{app:pcrit}):
\begin{enumerate}[leftmargin=2em]
  \item \textbf{Frobenius distinguishability:} $\|\Gam - I/7\|_F^2 > \|I/7\|_F^2$
        requires $P > 2/7$
  \item \textbf{Dominant eigenvalue:} $\lambda_{\max} \geq 2/7$ captures $\geq 2/7$
        of the total weight
  \item \textbf{Bayesian distinguishability:} expected posterior ratio
        $\mathbb{E}[\Lambda] > 1$ requires $P > 2/N$
  \item \textbf{Fano channel:} coherence contraction factor $1/3$ preserves structure
        only when $P > 2/7$
  \item \textbf{Self-observation:} minimum reflection $R \geq 1/3$ is self-consistent
        with $P \geq \Pcrit$
\end{enumerate}
\end{theorembox}

The convergence of five independent arguments on a single value is strong evidence that
$\Pcrit = 2/7$ is not an artifact of any single derivation method but a structural feature
of the geometry of $\Dens(\mathbb{C}^7)$.

\begin{figure}[H]
\centering
% Phase Diagram: Consciousness regimes in P-R plane
% ---------------------------------------------------------------
% IMPORTANT: R = 1/(7P) is a strict function of P (proved in §3.4
% from R := 1 - ||Gamma - I/7||_F^2 / ||Gamma||_F^2). Therefore the
% set of all physically realisable (P,R) pairs is the 1-dimensional
% curve R = 1/(7P), not a 2-D region. The colour bands below are
% intervals of P along that curve, drawn as vertical strips only to
% mark the corresponding P-range; they do NOT represent 2-D phases.
%
% Physical domain: P in [1/N, 1] = [1/7, 1]. Values P < 1/7 are
% forbidden by Cauchy-Schwarz for any density matrix on C^7 and are
% excluded from the plot.
% ---------------------------------------------------------------

\begin{tikzpicture}
\begin{axis}[
  width=12cm, height=9cm,
  xlabel={Purity $P = \mathrm{Tr}(\Gamma^2)$},
  ylabel={Reflection $R = 1/(7P)$},
  xmin=0.12, xmax=1.02,
  ymin=0.12, ymax=1.08,
  grid=both,
  grid style={gray!20},
  legend pos=north east,
  legend style={font=\scriptsize, fill=white, fill opacity=0.72,
                text opacity=1, draw=gray!50},
  title={\textbf{Consciousness phase diagram:} physical states
         lie on $R = 1/(7P)$, $P \in [1/7,\,1]$},
  title style={font=\small},
  every axis label/.style={font=\small},
  xtick={0.143,0.286,0.429,0.571,0.714,0.857,1.0},
  xticklabels={$\tfrac{1}{7}$,$\tfrac{2}{7}$,$\tfrac{3}{7}$,
               $\tfrac{4}{7}$,$\tfrac{5}{7}$,$\tfrac{6}{7}$,$1$},
  ytick={0.143,0.333,0.5,1.0},
  yticklabels={$\tfrac{1}{7}$,$\tfrac{1}{3}$,$\tfrac{1}{2}$,$1$},
]

% === Background strips marking P-intervals along the curve ===
% Layout strategy for band labels:
%  - The legend occupies the top-right corner (y > 0.92, x > 0.77).
%  - The physical curve R = 1/(7P) descends from (1/7, 1) to (1, 1/7).
%  - The I/7 thermal-death marker label lives at the extreme top-left
%    (~x<0.26, y~0.97), so the death-spiral band label is placed BELOW
%    the curve to avoid that label.
%  - The Goldilocks and Over-rigid band labels are placed ABOVE the curve
%    at a common height y = 0.85, which is clear of the legend bottom
%    (y~0.92) and of the curve (R<0.4 across those bands).
%  - Horizontally, the Over-rigid label is shifted LEFT within its wide
%    band (x=0.60) so that it does not extend into the legend column.

% Sub-threshold [1/7, 2/7): death spiral trajectories
\fill[red!12] (axis cs:0.143,0.12) rectangle (axis cs:0.286,1.08);
\node[font=\scriptsize, text width=1.5cm, align=center,
      color=red!60!black] at (axis cs:0.215, 0.22)
  {Death\\spiral};

% Goldilocks (2/7, 3/7]: all four L2 conditions simultaneously achievable
% Narrow band (~0.143 wide); use scriptsize bold with 2-line wrap so the
% label fits strictly inside the band.
\fill[green!15] (axis cs:0.286,0.12) rectangle (axis cs:0.429,1.08);
\node[font=\scriptsize\bfseries, color=green!45!black,
      text width=1.3cm, align=center]
  at (axis cs:0.358, 0.85)
  {Goldilocks\\(L2)};

% Over-rigid region (3/7, 1]: R<1/3 AND D_diff<2 (T-124)
% Shifted LEFT within the band to avoid the legend column (x>0.77).
\fill[gray!12] (axis cs:0.429,0.12) rectangle (axis cs:1.02,1.08);
\node[font=\scriptsize, color=gray!60!black,
      text width=2.4cm, align=center]
  at (axis cs:0.598, 0.85)
  {Over-rigid (T-124):\\$R<\tfrac{1}{3}$,\ $D_{\mathrm{diff}}<2$};

% === Threshold lines ===
\addplot[thick, red, dashed] coordinates {(0.286, 0.12) (0.286, 1.08)};
\node[font=\scriptsize, red, anchor=south east, rotate=90]
  at (axis cs:0.286, 0.14) {$P_{\mathrm{crit}} = 2/7$};

\addplot[thick, blue, dashed] coordinates {(0.12, 0.333) (1.02, 0.333)};
\node[font=\scriptsize, blue, anchor=west] at (axis cs:0.73, 0.305)
  {$R_{\mathrm{th}} = 1/3$};

\addplot[thick, green!50!black, dotted] coordinates {(0.429, 0.12) (0.429, 1.08)};
\node[font=\scriptsize, green!50!black, anchor=south east, rotate=90]
  at (axis cs:0.429, 0.14) {$P = 3/7$};

% === Physical curve R = 1/(7P) on full physical range [1/7, 1] ===
\addplot[very thick, plospurple, domain=0.143:1.0, samples=120] {1/(7*x)};
\addlegendentry{$R = 1/(7P)$ (physical states)}

% === Highlighted Goldilocks arc, half-open at P = 2/7 ===
\addplot[ultra thick, green!50!black, domain=0.2870:0.429, samples=40]
  {1/(7*x)};
\addlegendentry{Goldilocks arc, $P \in (2/7,\,3/7]$}

% === Key points: open at 2/7 (excluded), closed at 3/7 (included) ===
% (2/7, 1/2): open circle -- strict inequality P > 2/7
\addplot[only marks, mark=o, mark size=3pt, thick, black] coordinates {
  (0.286, 0.5)
};
% (3/7, 1/3): filled circle -- included in Goldilocks
\addplot[only marks, mark=*, mark size=3pt, black] coordinates {
  (0.429, 0.333)
};
% (1, 1/7): pure state endpoint
\addplot[only marks, mark=*, mark size=2pt, gray!60!black] coordinates {
  (1.0, 0.1429)
};

\node[font=\scriptsize, anchor=south west] at (axis cs:0.300, 0.52)
  {$(2/7,\,\tfrac{1}{2})$: viability (excluded)};
\node[font=\scriptsize, anchor=north west] at (axis cs:0.435, 0.315)
  {$(3/7,\,\tfrac{1}{3})$: upper L2 bound};
\node[font=\scriptsize, anchor=south, color=gray!60!black]
  at (axis cs:1.0, 0.16) {pure};

% Thermal-death endpoint at P = 1/7 (boundary of physical domain)
\addplot[only marks, mark=*, mark size=2.5pt, red!70!black] coordinates {
  (0.143, 1.0)
};
\node[font=\scriptsize, anchor=west, color=red!70!black]
  at (axis cs:0.152, 0.97) {$I/7$: thermal death};

\end{axis}
\end{tikzpicture}

\caption{\textbf{Phase diagram in the $P$--$R$ plane.} Because
$R = 1/(7P)$ is a strict functional identity (Eq.~\ref{eq:reflection}),
the set of physically realisable $(P,R)$ pairs is the \emph{one-dimensional}
purple curve, not a 2D region; the colour strips are intervals of $P$ along
that curve, not 2D phases. The physical domain is
$P \in [1/7,\,1]$: values $P < 1/7$ are forbidden by Cauchy--Schwarz for any
density matrix on $\mathbb{C}^7$, and the point $P = 1/7$ is the thermal
death state $I/7$ (red marker). The Goldilocks arc
$P \in (2/7,\,3/7]$ is the interval where all L2 conditions can be
satisfied simultaneously: the open circle at $(2/7,\,1/2)$ marks the
viability threshold (\emph{excluded}, strict inequality $P > 2/7$) and
the filled circle at $(3/7,\,1/3)$ marks the upper boundary
(\emph{included}). For $P > 3/7$ the L2 criterion fails by two
independent constraints: $R < 1/3$ \emph{and}
$D_{\mathrm{diff}} < D_{\min} = 2$ (T-124~\status{T}); the latter is
not visible in this 2D projection.}
\label{fig:phase-diagram}
\end{figure}

\paragraph{Comparison with empirical data.}
The Perturbational Complexity Index (PCI) was introduced by Casali et
al.~\cite{casali2013} on 32 healthy subjects (plus sleep and sedation
conditions). It was subsequently validated on a benchmark of 150 subjects
(healthy controls plus communicative brain-injured subjects in various states
of consciousness) by Casarotto et al.~\cite{casarotto2016}, who identified an
empirical cutoff $\mathrm{PCI}^* = 0.31$ that discriminated conscious from
unconscious states with 100\,\% sensitivity and specificity on the benchmark,
and 94.7\,\% sensitivity in a clinical MCS cohort. The theoretical prediction
$\Pcrit = 2/7 \approx 0.286$ differs from the empirical cutoff by approximately
$8\,\%$, within the calibration uncertainty of any neural-to-$\Gam$ mapping
$\pi_{\mathrm{bio}}$ (see Section~\ref{sec:experimental}). The value $2/7$ was
\emph{derived} from the theory, not fitted.

It must be emphasised that PCI and purity $P$ are mathematically distinct
quantities: PCI is a normalised Lempel--Ziv complexity of TMS-evoked
spatiotemporal responses, while $P = \Tr(\Gam^2)$ is a Hilbert--Schmidt norm
on density matrices. The numerical proximity $0.286 \approx 0.31$ is
suggestive but does not itself constitute evidence; it becomes evidence only
once the mapping $\pi_{\mathrm{bio}}$ is calibrated (Phase~II,
\S\ref{sec:experimental}). The coincidence is presented as motivation for
experimental investigation, not as empirical confirmation.

\subsection{The evolution equation}
\label{subsec:evolution}

The dynamics of $\Gam$ are governed by the Liouvillian superoperator $\Lind_\Omega$,
derived from axioms A1--A4 via the Page--Wootters constraint and the structure of $\Omega$
(Figure~\ref{fig:evolution-flow}):

\begin{figure}[H]
\centering
% Evolution Equation: Three forces acting on Γ
% Clean layout — no overlapping labels

\begin{tikzpicture}[
  gammanode/.style={circle, draw=plospurple, fill=plospurple!15,
                    minimum size=2cm, align=center, font=\large\bfseries},
  forcebox/.style={rectangle, rounded corners=4pt,
                   minimum width=3cm, minimum height=1.4cm, align=center, font=\small},
  arw/.style={-{Stealth[length=3mm]}, very thick, #1},
  scale=1.0
]

% === Central Γ ===
\node[gammanode] (gamma) at (0, 0) {$\Gamma(\tau)$};

% === Three forces — spread wide ===
\node[forcebox, draw=blue!70, fill=blue!8] (unitary) at (-5, 3.5)
  {\textbf{Unitary}\\[1pt] $-i[H,\Gamma]$};

\node[forcebox, draw=red!70, fill=red!8] (dissipation) at (5, 3.5)
  {\textbf{Dissipation}\\[1pt] $\mathcal{D}_\Omega[\Gamma]$};

\node[forcebox, draw=green!60!black, fill=green!8] (regeneration) at (0, -4.5)
  {\textbf{Regeneration}\\[1pt] $\mathcal{R}[\Gamma, E]$};

% === Arrows with labels centered along each path ===
\draw[arw=blue!70] (unitary.south east) -- (gamma.north west)
  node[midway, above, sloped, font=\footnotesize, blue!70] {Preserves $P$};

\draw[arw=red!70] (dissipation.south west) -- (gamma.north east)
  node[midway, above, sloped, font=\footnotesize, red!70] {Decreases $P$};

\draw[arw=green!60!black] (regeneration.north) -- (gamma.south)
  node[midway, right=3pt, font=\footnotesize, green!50!black] {Increases $P$};

% === Compact descriptors under each box ===
\node[font=\scriptsize, gray, text width=3cm, align=center] at (-5, 2.2)
  {Energy-conserving\\rotation};
\node[font=\scriptsize, gray, text width=3cm, align=center] at (5, 2.2)
  {Decoherence\\(toward $I/7$)};
\node[font=\scriptsize, gray, text width=3.8cm, align=center] at (0, -6)
  {$\kappa = \kappa_{\mathrm{boot}} + \kappa_0 \cdot \mathrm{Coh}_E$\\[1pt]
   Strength $\propto$ quality of experience};

% === Balance equation — well below diagram ===
\node[font=\small, text width=9cm, align=center] at (0, -7.8) {
  $\displaystyle\frac{d\Gamma}{d\tau}
  = \underbrace{-i[H, \Gamma]}_{\Delta P = 0}
  + \underbrace{\mathcal{D}_\Omega[\Gamma]}_{\Delta P < 0}
  + \underbrace{\mathcal{R}[\Gamma, E]}_{\Delta P > 0}$
};

\node[font=\footnotesize\bfseries] at (0, -9.2) {
  \textcolor{green!50!black}{Life}: regeneration $\geq$ dissipation.
  \quad
  \textcolor{red!70}{Death}: regeneration $<$ dissipation.
};

\end{tikzpicture}

\caption{\textbf{The three forces governing the evolution of $\Gam$.} Unitary dynamics
(blue) preserves purity; dissipation (red) decreases it; regeneration (green) increases
it proportionally to E-coherence. Life is the regime where regeneration compensates
dissipation. Death is the regime where it does not.}
\label{fig:evolution-flow}
\end{figure}

\begin{equation}
  \frac{d\Gam(\tau)}{d\tau} = \Lind_\Omega[\Gam(\tau)]
  = \underbrace{-i[H, \Gam]}_{\text{unitary}}
  + \underbrace{\mathcal{D}_\Omega[\Gam]}_{\text{dissipation}}
  + \underbrace{\Regen[\Gam, E]}_{\text{regeneration}}
  \label{eq:evolution}
\end{equation}

\paragraph{Term 1: Unitary dynamics.}
The effective Hamiltonian emerges from the Page--Wootters mechanism (T-87~\status{T}):
\begin{equation}
  H_{\mathrm{eff}}(\tau) = H_{6D} + \langle\tau|H_{\mathrm{int}}|\tau\rangle_O
  \label{eq:heff}
\end{equation}
where $H_{6D}$ governs the six non-temporal dimensions and the second term arises from
tracing out the O-dimension (internal clock). Time $\tau$ is not an external parameter
--- it emerges from the correlation between $\Gam$ and the O-sector
(Section~\ref{subsec:emergent-time}).

\paragraph{Term 2: Dissipation (Lindblad).}
\begin{equation}
  \mathcal{D}_\Omega[\Gam] = \sum_k \gamma_k
  \left( L_k\, \Gam\, L_k^\dagger - \frac{1}{2}\{L_k^\dagger L_k,\, \Gam\} \right)
  \label{eq:dissipation}
\end{equation}

The Lindblad operators $L_k$ are derived from the atoms of the subobject classifier
$\Omega$ (L-unification theorem~\status{T}):
\begin{equation}
  L_k^{\mathrm{atom}} = |k\rangle\langle k|, \qquad k \in \{A,S,D,L,E,O,U\}
  \label{eq:lindblad-atoms}
\end{equation}

The Fano channel extends these to triple projectors (one per Fano line):
\begin{equation}
  L_p^{\mathrm{Fano}} = \frac{1}{\sqrt{3}}\,\Pi_p
  = \frac{1}{\sqrt{3}}\bigl(|i\rangle\langle i| + |j\rangle\langle j| + |k\rangle\langle k|\bigr)
  \label{eq:fano-operators}
\end{equation}
where $(i,j,k)$ is the $p$-th Fano line. Key properties of the Fano channel
$\mathcal{P}_{\mathrm{Fano}}$~\status{T}:
\begin{itemize}[leftmargin=2em]
  \item Diagonal preservation: $[\mathcal{P}_{\mathrm{Fano}}(\Gam)]_{ii} = \gamma_{ii}$
  \item Coherence contraction: $[\mathcal{P}_{\mathrm{Fano}}(\Gam)]_{ij} = \frac{1}{3}\gamma_{ij}$
        for $i \neq j$
  \item Phase preservation: $\arg([\mathcal{P}_{\mathrm{Fano}}(\Gam)]_{ij}) = \arg(\gamma_{ij})$
  \item Completeness: $\sum_{p=1}^{7} \Pi_p = 3I$ (each dimension on exactly 3 Fano lines)
\end{itemize}

Dissipation always decreases purity: $dP/d\tau|_{\mathcal{D}} \leq 0$~\status{T}. This
follows from the CPTP property of the Lindblad channel: CPTP maps are contractive in the
Bures metric (a consequence of the data-processing inequality), moving any state toward
thermal equilibrium $I/7$. Formally, for $\Gam \neq I/7$:
\begin{equation}
  \frac{dP}{d\tau}\bigg|_{\mathcal{D}} = -\Delta(\Lind_0) \cdot (P - 1/7) < 0
  \label{eq:dissipation-decreases-purity}
\end{equation}
where $\Delta(\Lind_0) > 0$ is the spectral gap of the linear Lindbladian (primitivity
theorem T-39a~\status{T}). Without regeneration, the system converges exponentially to
$I/7$ (thermal death).

\paragraph{Term 3: Regeneration.}
The regeneration channel $\Regen$ is the only mechanism that can increase purity. Its
strength depends on E-coherence:
\begin{equation}
  \kappa(\Gam) \;=\; \kappa_{\mathrm{boot}} \;+\; \kappa_0\,\CohE(\Gam).
  \label{eq:kappa}
\end{equation}
Here $\CohE(\Gam) = \|\pi_E(\Gam)\|_{\mathrm{HS}}^2 / \|\Gam\|_{\mathrm{HS}}^2$ is the
Hilbert--Schmidt projection of $\Gam$ onto the E-sector~\status{T}. The basal floor
\begin{equation}
  \kappa_{\mathrm{boot}} \;=\; \frac{\omega_0}{N} \;=\; \frac{\omega_0}{7}\qquad\status{O}
  \label{eq:kappa-boot}
\end{equation}
is the innate (``bootstrap'') regeneration rate set by the convention that the
minimal-symmetry unit of the adjunction $\mathcal{D}\dashv\Regen$ delivers one ``quantum''
of regeneration per dimension at scale $\omega_0$. The adaptive amplification
\begin{equation}
  \kappa_0(\Gam) \;=\; \omega_0\,\frac{|\gamma_{OE}|\,|\gamma_{OU}|}{\gamma_{OO}}
  \qquad\status{T}
  \label{eq:kappa-0}
\end{equation}
is the categorical norm $\|\mathrm{Nat}(\mathcal{D}_\Omega,\Regen)\|$ of the natural
transformation between the dissipation and regeneration functors (T-88~\status{T}, Yoneda
$+$ Bures $+$ Stinespring; see the master definition at
\url{https://holon.sh/docs/core/foundations/axiom-septicity#структурный-анзац-kappa0}).
Both $\kappa_{\mathrm{boot}}$ and $\kappa_0$ are \emph{derived}, not fitted. The full
regeneration term is
\begin{equation}
  \Regen[\Gam, E] \;=\; \kappa(\Gam) \,\bigl(\rho^* - \Gam\bigr)\, g_V(P),
  \label{eq:regen-full}
\end{equation}
where $\rho^* = \vp(\Gam)$ is the categorical self-model (T-96~\status{T}) and $g_V(P)$
is the viability gate: $g_V(P) = 1$ for $P > \Pcrit$, and $g_V(P) = 0$ otherwise
(Landauer limit --- below threshold, free energy $\Delta F \leq 0$ forbids regeneration).
The proportionality $\kappa \propto \CohE$ is the mathematical content of the claim that
\emph{consciousness drives survival}: the quality of inner experience directly
determines the rate of regeneration.

\paragraph{Thermodynamic interpretation.}
The viability gate $g_V(P)$ has a thermodynamic origin. Regeneration requires free energy
$\Delta F > 0$ --- the system must do work against entropy to increase purity. At $P = \Pcrit$,
$\Delta F = 0$ exactly: the system is at thermodynamic equilibrium with its dissipative
environment. Below $\Pcrit$, $\Delta F < 0$ (Landauer's principle~\cite{landauer1961}),
and regeneration is thermodynamically forbidden --- not merely insufficient, but impossible.
Free energy $\Delta F > 0$ thus plays the role of Schr\"{o}dinger's ``negentropy''
\cite{schrodinger1944}: it is the fuel for regeneration of purity. The death spiral
(Eq.~\ref{eq:death-spiral}) is not a failure of mechanism but a thermodynamic
inevitability: below threshold, the second law guarantees irreversible decay.

\subsection{The self-modeling operator $\vp$}
\label{subsec:phi-operator}

\begin{definition}[Self-modeling operator]
The self-modeling operator is a CPTP (Completely Positive Trace-Preserving) channel:
\begin{equation}
  \vp: \Dens(\Hil) \to \Dens(\Hil), \qquad
  \vp(\Gam) = \sum_m K_m\, \Gam\, K_m^\dagger, \qquad
  \sum_m K_m^\dagger K_m = I
  \label{eq:phi-operator}
\end{equation}
\end{definition}

The system's self-model $\vp(\Gam)$ is an \emph{approximate} copy of $\Gam$, not an exact
one (the no-cloning theorem forbids exact copying in quantum mechanics). The canonical
physical realization (T-62~\status{T}):
\begin{equation}
  \vp_k(\Gam) = (1-k)\Gam + k\rho^*
  \label{eq:phi-physical}
\end{equation}
where $\rho^* = \vp(\Gam)$ is the categorical self-model (the fixed point of the
self-observation process, T-96~\status{T}) and $k = 1 - R$ is the compression parameter.

This definition is not circular: $\rho^*$ is constructed independently via the categorical left adjoint (T-96~\status{T}), yielding a fixed point that $\vp$ converges to iteratively. Starting from any initial guess $\rho^*_0 = I/7$, the sequence $\rho^*_{n+1} = \vp(\Gam; \rho^*_n)$ converges exponentially (Banach fixed-point theorem, Eq.~\ref{eq:fixed-point}).

\paragraph{Fixed-point theorem \status{T} (Banach).}
Iterating $\vp$ converges to a fixed point:
\begin{equation}
  \vp^n(\Gam_0) \to \Gam^*, \qquad
  \|\vp^n(\Gam_0) - \Gam^*\|_F \leq k^n \|\Gam_0 - \Gam^*\|_F
  \label{eq:fixed-point}
\end{equation}
This convergence is the mathematical content of ``stable self-knowledge'': the system's
self-model approaches a fixed point that represents its best possible self-understanding.

\subsection{The reflection measure $R$}
\label{subsec:reflection}

\begin{definition}[Reflection measure]
\begin{equation}
  R(\Gam) := 1 - \frac{\|\Gam - I/7\|_F^2}{\|\Gam\|_F^2} = \frac{1}{7P(\Gam)}
  \label{eq:reflection}
\end{equation}
\end{definition}

$R$ measures how well the system models itself, ranging from $R = 1$ (perfect self-model,
only at $P = 1/7$ --- trivially, when there is nothing to model) to $R \to 0$ as $P \to \infty$
(which is unphysical since $P \leq 1$). In the physical range:

\begin{itemize}[leftmargin=2em]
  \item $R = 1$: at $P = 1/7$ (thermal death) --- trivially, the system ``knows itself''
        because there is nothing to know (uniform noise matches any model)
  \item $R = 1/3$: at $P = 3/7$ --- upper boundary of the consciousness window
        ($R$ drops below $\Rth$ for $P > 3/7$)
  \item $R = 1/(7 \cdot 2/7) = 1/2$: at $P = 2/7$ (viability threshold)
\end{itemize}

The inverse relationship between $R$ and $P$ reflects a fundamental feature: a more organized system (higher $P$) is further from its dissipative attractor $I/7$, making self-modeling \emph{harder} (lower $R$). The L2 consciousness window $P \in (2/7, 3/7]$ is thus a Goldilocks zone: enough structure to be viable, but not so much that self-modeling becomes impossible. Above $P = 3/7$, the system is too rigid for adaptive self-observation --- differentiation $D_{\mathrm{diff}}$ drops below $\Dmin = 2$ (T-124~\status{T}).

\begin{theorembox}[Reflection threshold (T-40b) \status{T}]
The reflection threshold for cognitive (L2) consciousness is:
\begin{equation}
  \Rth = \frac{1}{3} \approx 0.333
  \label{eq:rth}
\end{equation}
derived from the triadic decomposition $K = 3$ of Lindblad generators (T-40a~\status{T},
T-57~\status{T}) combined with Bayesian dominance of the self-model over the $K$
alternatives. Below $R = 1/3$, the system exists but does not model itself sufficiently
for reflective awareness.
\end{theorembox}

\subsection{Integration and differentiation}
\label{subsec:integration}

\begin{definition}[Integration measure $\Phi$]
\begin{equation}
  \Phi(\Gam) = \frac{\sum_{i \neq j} |\gamma_{ij}|^2}{\sum_i \gamma_{ii}^2}
  = \frac{P_{\mathrm{off}}}{P_{\mathrm{diag}}}
  \label{eq:phi}
\end{equation}
\end{definition}

$\Phi$ measures the ratio of off-diagonal coherence (binding between dimensions) to
diagonal weight (individual dimension strength). When $\Phi \geq 1$, the system is
\emph{coherence-dominated}: the connections between dimensions are at least as strong
as the dimensions themselves. This is the mathematical content of ``unity of experience.''

\begin{theorembox}[Integration threshold (T-129) \status{T}]
The integration threshold for L2 consciousness is:
\begin{equation}
  \Phith = 1
  \label{eq:phith}
\end{equation}
This is the unique self-consistent value at $\Pcrit = 2/7$: coherent dominance
($\Phi \geq 1$) requires that off-diagonal coherences collectively outweigh diagonal
elements.
\end{theorembox}

\begin{definition}[Experience differentiation $D_{\mathrm{diff}}$]
\begin{equation}
  D_{\mathrm{diff}}(\Gam) = \exp\bigl(S_{\mathrm{vN}}(\rho_E)\bigr)
  \label{eq:ddiff}
\end{equation}
where $S_{\mathrm{vN}}(\rho_E) = -\Tr(\rho_E \ln \rho_E)$ is the von Neumann entropy
of the E-sector reduced density matrix.
\end{definition}

$D_{\mathrm{diff}}$ measures the diversity of accessible experience states. The threshold
$\Dmin = 2$~\status{T} (T-151) means the system must distinguish at least two qualitatively
different experiences for consciousness to obtain.\footnote{The partial trace
$\rho_E = \Tr_{-E}(\Gam)$ requires the extended 42D formalism (since 7 is prime and
$\Hil$ does not factor as a tensor product). In the minimal 7D setting, $D_{\mathrm{diff}}$
is approximated via the spectrum of $\Gam$ restricted to the E-sector. See
\url{https://holon.sh/docs/core/structure/dimension-e} for the full treatment.}

\subsection{Consciousness measure $C$}
\label{subsec:consciousness-measure}

\begin{definition}[Consciousness measure]
\begin{equation}
  C(\Gam) := \Phi(\Gam) \times R(\Gam)
  \label{eq:consciousness-measure}
\end{equation}
\end{definition}

$C$ quantifies the degree of conscious experience as the product of \emph{integration}
(how much dimensions are bound into a whole) and \emph{reflection} (how well the system
models itself). Why the product, not the sum? Because integration without reflection means
bound but unconscious (like a coma patient whose brain may show residual integration);
reflection without integration means self-aware but fragmented (like dissociative states).
Only the product captures ``integrated self-awareness'' --- the conjunction that defines
L2 consciousness.

\subsection{Levels of consciousness}
\label{subsec:levels}

\begin{figure}[H]
\centering
% Consciousness Levels Hierarchy: L0 → L1 → L2 → L3

\begin{tikzpicture}[
  levelbox/.style={rectangle, rounded corners=4pt,
                minimum width=10cm, minimum height=1.4cm, align=center, font=\small},
  threshold/.style={font=\scriptsize\bfseries, #1},
  arrow/.style={-{Stealth[length=2.5mm]}, thick, gray},
  scale=1.0
]

% === Levels (bottom to top) ===
\node[levelbox, draw=gray, fill=gray!15]         (l0)   at (0, 0)   {\textbf{L0: Formal interiority} ($\Gamma \in \mathcal{D}(\mathcal{H})$)\\
  \footnotesize Any system has an internal aspect. Examples: electron, stone};

\node[levelbox, draw=orange, fill=orange!15]    (l1)   at (0, 2.2) {\textbf{L1: Proto-consciousness} ($\mathrm{rank}(\rho_E) > 1$)\\
  \footnotesize Non-trivial E-sector. Examples: thermostat, bacterium};

\node[levelbox, draw=blue, fill=blue!15]      (l2)   at (0, 4.4) {\textbf{L2: Cognitive consciousness}
  ($P > 2/7,\; R \geq 1/3,\; \Phi \geq 1,\; D \geq 2$)\\
  \footnotesize Full qualia, self-awareness. Examples: waking humans, corvids, primates};

\node[levelbox, draw=green!50!black, fill=green!15] (l3) at (0, 6.6) {\textbf{L3: Meta-consciousness}
  (L2 + higher-order $R^{(n)}$)\\
  \footnotesize Reflects on reflecting. Examples: meditation, philosophical inquiry};

\node[levelbox, draw=plospurple, fill=plospurple!15]   (l4)   at (0, 8.8) {\textbf{L4: Limit}
  (SAD $= 3$, T-142)\\
  \footnotesize Maximum self-reflection depth. Theoretical ceiling};

% === Threshold annotations (between levels) ===
\draw[arrow] (l0.north) -- (l1.south);
\node[threshold=orange, fill=white, inner sep=2pt] at (0, 1.1)
  {$\mathrm{rank}(\rho_E) > 1$ \quad [non-trivial E-sector]};

\draw[arrow] (l1.north) -- (l2.south);
\node[threshold=blue, fill=white, inner sep=2pt] at (0, 3.3)
  {$P > 2/7,\; R \geq 1/3,\; \Phi \geq 1,\; D \geq 2$ \quad [T-129, T-151]};

\draw[arrow] (l2.north) -- (l3.south);
\node[threshold=green!50!black, fill=white, inner sep=2pt] at (0, 5.5)
  {Higher-order $R^{(n)} \geq R_{\mathrm{th}}^{(n)}$};

\draw[arrow] (l3.north) -- (l4.south);
\node[threshold=plospurple, fill=white, inner sep=2pt] at (0, 7.7)
  {SAD$_{\max} = 3$ [T-142] --- geometric ceiling};

\end{tikzpicture}

\caption{\textbf{Hierarchy of consciousness levels L0--L4.} Each transition requires
crossing a mathematical threshold. L0 is universal (any $\Gam$); L2 requires all four
conditions simultaneously; L4 is the geometric ceiling imposed by Fano contraction.}
\label{fig:levels}
\end{figure}

\begin{table}[H]
\centering
\caption{Levels of consciousness and their thresholds. All thresholds are theorems, not postulates.}
\label{tab:levels}
\small
\begin{tabular}{@{}clp{4.5cm}p{3.8cm}l@{}}
\toprule
\textbf{Level} & \textbf{Name} & \textbf{Conditions} & \textbf{What it means} & \textbf{Example} \\
\midrule
L0 & Formal interiority & $\Gam \in \Dens(\Hil)$ & Any system has an internal aspect & Electron \\
L1 & Proto-consciousness & $\mathrm{rank}(\rho_E) > 1$ & Non-trivial E-sector & Thermostat \\
L2 & Cognitive & $P > \tfrac{2}{7},\; R \geq \tfrac{1}{3},\; \Phi \geq 1,\; D_{\mathrm{diff}} \geq 2$ & Full qualia, self-awareness & Waking human \\
L3 & Meta & L2 $+\; R^{(2)} \geq 1/4$ & Reflects on reflecting & Meditation \\
L4 & Limit & $\mathrm{SAD} = 3$ (T-142) & Maximum self-reflection depth & Theoretical ceiling \\
\bottomrule
\end{tabular}
\end{table}

\subsection{Emergent time}
\label{subsec:emergent-time}

Time in UHM is not postulated as an external parameter. It emerges from the internal
structure of $\Gam$ via the Page--Wootters mechanism (T-87~\status{T}).

The O-dimension (ground/energy) acts as an internal clock. The total system satisfies the
constraint:
\begin{equation}
  [\hat{C},\, \Gam_{\mathrm{total}}] = 0
  \label{eq:pw-constraint}
\end{equation}
where $\hat{C} = H_O \otimes \mathbb{1}_{6D} + \mathbb{1}_O \otimes H_{6D} + H_{\mathrm{int}}$
is the total energy operator. Conditioning on the O-sector yields effective dynamics:
\begin{equation}
  i\frac{\partial}{\partial\tau}\Gam(\tau) = [H_{\mathrm{eff}}(\tau),\, \Gam(\tau)]
  \label{eq:pw-dynamics}
\end{equation}

Discrete time from the temporal modality $\triangleright: S_i \mapsto S_{(i+1) \bmod 7}$
gives $\tau_n = \triangleright^n(\mathrm{now})$ with period $\mathbb{Z}_7$. In the
thermodynamic limit of $M$ coupled holons: $\mathbb{Z}_{7^M} \to \mathbb{R}$
(continuous time emerges, T-118~\status{T}).
